
\subsection*{BTMs: Bayesian hierarchical models}

The SBM structure was built around the data collection process, with studies that sampled sites which were sometimes organized into spatial blocks. In the BTMs, we instead created the hierarchical structure based on the biome and realm where each site is located, and the taxonomic group(s) that were sampled at that site. This was implemented as a two-level hierarchy. At the first level, biomes were combined with taxonomic groupings. Biomes (e.g. tropical and subtropical moist broadleaf forests) are distinct in their environmental conditions and underlying biodiversity patterns, which suggests that the effect sizes of both natural and anthropogenic drivers will be different between them. The taxonomic groupings were based on differences in spatial distribution patterns\cite{jenkins2015} and potentially differentiated responses to drivers. At the same time, data limitations and taxonomic imbalances implied that most taxa were aggregated at quite high levels. For the second hierarchical level, the biome-taxonomic groups were further subdivided by biogeographic realms (e.g. Neotropics) to form so-called regional biomes\cite{bevan2024}. This results in a structure characterized by parent-child relationship between groups at the first and second hierarchical levels.

Although this model could have been implemented as a frequentist linear mixed model, the Bayesian formulation enabled additional flexibility, stronger borrowing of statistical strength across groups, and regularization through weakly informative priors \cite{dorazio2016,gelman2012,lemoine2019}. Groups with low information content are pulled towards their parent group, while groups with sufficient statistical power are able to deviate from those mean values. Still, note that the main comparison in this study is between the overall model approach and structure, not frequentist versus Bayesian implementations.

We let $b = 1, \dots, B$ denote biomes, $t = 1, \dots, T$ taxonomic groups, and $r = 1, \dots, R$ biogeographic realms. Further, $r(b)$ indicates the partition of a given biome into realms. However, note that realms can span multiple biomes and vice versa, so this is not a unique one-to-one mapping. For the group-level model parameters, intercepts are denoted by $\alpha$ and other parameters by $\boldsymbol{\beta}$. To make observations conditionally independent during training, we also included study-block identifiers $\gamma_{sb(s)}$ as control variables, identical to the random intercepts of the SBM models. These were used in model training but not for any out-of-sample predictions. Like the SBMs, the BTMs assumed a beta distributed likelihood with a logit link function. The alpha diversity (GMA) model for taxonomic group $t$ at site $i$, in the regional biome $r(b)$ can be written as
\vspace{0.25\baselineskip}
\begin{equation*}
    g(\mu_{i,t}) = \alpha_{t,r(b)} + \mathbf{x}_i^\top \boldsymbol{\beta}_{t,r(b)} + \gamma_{sb(s)},
\end{equation*}

where the covariates $\mathbf{x}_i^\top$ are defined at the site level (same for all taxonomic units $t$ at the site). We assumed normal and uncorrelated priors for the model parameters, with half-normal priors on the variance terms. At the population-level, the following hyperpriors were used:
\vspace{0.25\baselineskip}
\begin{equation*}
\begin{aligned}
    \mu_{\alpha} \sim \mathcal{N}(0.35, 0.30^2), \; \mu_{\beta}^{(k)} \sim \mathcal{N}(0, 0.15^2)  \; \text{for} \; k = 1, \dots, K,
\end{aligned}
\end{equation*}

where the $K$ is the number of regression parameters. These priors define distributions over the population-level model parameters, which sit at the very top of the hierarchy. The priors at the first hierarchical level (biome-taxa) were drawn from the hyperpriors in the following way:
\vspace{0.25\baselineskip}
\begin{equation*}
\begin{gathered}
    \alpha: \sigma_{\alpha} \sim \text{Half-Normal}(0.20^2), \; \alpha_{t,b} \sim \mathcal{N}(\mu_{\alpha}, \tau_{t,b} \cdot \sigma_{\alpha}^2) \\
    \beta: \sigma_{\beta} \sim \text{Half-Normal}(0.12^2), \; \beta_{t,b}^{(k)} \sim \mathcal{N}(\mu_{\beta}, \tau_{t,b} \cdot \sigma_{\beta}^2)
\end{gathered}
\end{equation*}

The scaling factor, $\tau_{t,b} = \sqrt{n_{t,b}} -1$, adapted the amount of regularization -- through the prior variance -- on the group-level model parameters as a function of the number of studies in each biome-taxa group. In a similar way, the priors at the second level (biome-taxa-realm) inherited from their respective parent groups at the first level, but with tighter priors on the slope variances:
\vspace{0.25\baselineskip}
\begin{equation*}
\begin{gathered}
    \alpha: \sigma_{\alpha} \sim \text{Half-Normal}(0.20^2), \; \alpha_{t,r(b)} \sim \mathcal{N}(\alpha_{t,b}, \tau_{t,r(b)} \cdot \sigma_{\alpha}^2), \\
    \beta: \sigma_{\beta} \sim \text{Half-Normal}(0.10^2), \; \beta_{t,r(b)}^{(k)} \sim \mathcal{N}(\beta_{t,b}^{(k)}, \tau_{t,r(b)} \cdot \sigma_{\beta}^2).
\end{gathered}
\end{equation*}

The numeric prior values above were chosen based iterative prior predictive checks, to obtain a prior predictive distribution reasonably similar to the observed data (\extfigpanel{fig:response_data_distr}{c}) and achieve a relatively high degree of model regularization. Next, the beta diversity model can be expressed as
\vspace{0.25\baselineskip}
\begin{equation*}
    g(\mu_{ij,t}) = \alpha_{t,r(b)} + \mathbf{x}_i^\top \boldsymbol{\beta}_{t,r(b)} +\mathbf{z}_{ij}^\top \boldsymbol{\beta}_{t,r(b)}^\Delta + \gamma_{sb(s)}.
\end{equation*}

The prior structure was identical to the alpha diversity model, but we assumed slightly different hyperpriors due to the lower mean and longer tail of the BC similarity distribution (\extfigpanel{fig:response_data_distr}{d}):
\vspace{0.25\baselineskip}
\begin{equation*}
\begin{aligned}
    \mu_{\alpha} \sim \mathcal{N}(0.25, 0.25^2), \; \sigma_{\alpha} \sim \text{Half-Normal}(0.15^2).
\end{aligned}
\end{equation*}

Further, the scaling factor used the natural logarithm, $\tau_{t,b} = \text{ln}(n_{t,b}) -1$, since this gave better results. For the dispersion parameter of the Beta likelihood, we sampled a raw scale
parameter $\sigma_{\mathrm{raw}} \sim \mathrm{Beta}(a,b)$, using $a=2,b=12$ for
the GMA model and $a=2,b=20$ for the BC similarity model. We then defined a mean-dependent dispersion
$\sigma(\mu_{i,t}) = \sigma_{\mathrm{raw}} \sqrt{\mu_{i,t}\left(1-\mu_{i,t}\right)}$ and used it to parameterize the beta likelihood:
\begin{equation*}
    y_{i,t} \sim \mathrm{Beta}\!\left(\mu_{i,t}, \sigma(\mu_{i,t})\right).
\end{equation*}

Since many biogeographic and taxonomic groups contained only a few studies and sites (\autoref{fig:predicts_data}), there was a risk that the estimated group-level parameters would become too specific to a few studies, rather than generally applicable to the group at large. To prevent this from producing overoptimistic accuracy numbers, especially for interpolation, we implemented a 'roll-up' scheme for out-of-sample predictions. Here, biome-taxa-realm ($t,r(b)$) groups with less than five studies were assigned the parameters of their respective parent groups at the biome-taxa ($t,b$) level for prediction, such that $\boldsymbol{\beta}_{t,r(b)} := \boldsymbol{\beta}_{t,r}$. If the threshold was still not met at the biome-taxa level, the population-level posterior parameters were used. The BTMs were implemented in \texttt{PyMC} \cite{abril-pla2023}, using the No-U-Turn Sampler (NUTS)\cite{hoffman2014} with a \texttt{numpyro}\cite{phan2019} backend. For all experiments, we ran four parallel sampling chains for stability, using 1,000 tuning samples that were discarded and 1,000 posterior draws. A high target acceptance rate of 0.95 was used to avoid problems when sampling from the complex posterior distribution. Due to the high number of diverse hierarchical groups before roll-up (181), some which are quite small, it was hard to avoid some parameters having small effective sample sizes and $\widehat{R}>1.01$\cite{vehtari2019}. While this could indicate convergence issues in the sampling, the overall performance results were robust also when using much fewer iterations than in the final runs.

\vspace{\baselineskip}
